% !TeX root = ../main.tex
Pooled CRISPR screens increasingly measure guide abundance across more than
two experimental conditions.  Examples include time courses, dose series,
ordered FACS fractions, donor-stratified experiments, and factorial designs.
In these studies the scientific estimand is a trend, an adjusted association,
or an interaction.  Reducing the experiment to a high-versus-low comparison
discards intermediate observations and prevents direct adjustment for
additional design variables.

The sampling model is consequential in this setting.  RNA-sequencing methods
typically model a transcript count with a negative-binomial distribution and
represent library size through normalization or an offset
\citep{anders2010deseq,love2014deseq2,law2014voom,lun2016edger}.  Several
CRISPR-screen methods adopted this framework
\citep{li2014mageck,li2015mageckvispr}.  A pooled-guide count also has a known
denominator, however: the total number of mapped guide reads in that library.
Conditional on this total, guide abundance is a proportion.  The
beta-binomial model represents this structure directly and separates
depth-dependent sampling variation from heterogeneity among biological
libraries.  In the original CB\textsuperscript{2} study, goodness-of-fit
comparisons favored the beta-binomial over negative-binomial alternatives for
pooled CRISPR counts \citep{jeong2019beta}.

Beta-binomial regression itself is not new.  The \texttt{corncob} framework
models both the mean and overdispersion of counts conditional on library totals
\citep{martin2020corncob}, and general implementations are available in
\texttt{aod}, \texttt{VGAM}, and \texttt{gamlss}
\citep{lesnoff2012aod,yee1996vgam,rigby2005gamlss}.  BARCS is therefore not
presented as the first beta-binomial regression method.  Its contribution is a
CRISPR-screen implementation that preserves immutable guide-library totals,
uses an R design matrix, returns guide-level coefficient tests at screen scale,
and connects this workflow to CB\textsuperscript{2}.

CB\textsuperscript{2} applies its sampling model to two independent groups.
For guide $g$ observed in library $i$, let $K_{gi}$
denote the guide count, $N_i$ the unfiltered mapped-guide total, and
$\mathbf{x}_i^\top$ one row of a prespecified, full-rank sample design matrix.
The regression problem is
\begin{equation}
  K_{gi}\mid N_i
  \sim \operatorname{BetaBinomial}(N_i,\mu_{gi},\rho_g),
  \qquad
  \logit(\mu_{gi})=\mathbf{x}_i^\top\boldsymbol{\beta}_g,
  \qquad
  H_0:\mathbf{c}^\top\boldsymbol{\beta}_g=0 .
  \label{eq:intro-target}
\end{equation}
Here time, dose, donor, batch, or treatment is a sample-level predictor; the
guide proportion remains the response.  Independent biological libraries,
not individual reads, determine the residual degrees of freedom.

Existing methods address parts of this problem.  Repeated two-group analyses
cannot estimate one coefficient conditional on the remaining columns of the
design matrix.  General RNA-sequencing tools support arbitrary designs but do
not condition each guide on its observed library total.  At the other end of
the spectrum, Chronos and Waterbear model population dynamics or joint FACS-bin
membership, respectively, but are tied to those experimental structures
\citep{dempster2021chronos,pimentel2024waterbear}.  A general-purpose,
depth-aware coefficient test for pooled CRISPR designs remains missing.

We introduce BARCS (Beta-binomial Analysis and Regression for CRISPR Screens)
to fill this gap.  BARCS retains the beta-binomial sampling principle of
CB\textsuperscript{2}, replaces the two group means with an R model matrix, and
tests a prespecified coefficient or contrast with sample-level inference.  The
relationship to CB\textsuperscript{2} is structural: both methods condition
guide counts on library totals and cap the information contributed by deeper
sequencing.  Their estimators, effect scales, dispersion assumptions, and
finite-sample reference distributions differ, so BARCS is not claimed to nest
the implemented CB\textsuperscript{2} test.

The empirical evaluation begins with calibration under known truth.  It then
uses the three-time-point Cas13 study of four replicate-complete cell lines from
\citet{liang2026lncrna} as a processed-count sensitivity analysis and an
ordered-bin IL2RA screen as a held-out-control diagnostic.  Genome-scale and
factorial simulations examine dispersion moderation and denominator choice.
The serially harvested HT-29 culture, endpoint screens, and copy-number
analyses are retained as descriptive robustness checks rather than evidence for
sample-level longitudinal inference.

This design supports a deliberately limited claim.  BARCS makes a
library-total-conditional beta-binomial mean model available to designs that
CB\textsuperscript{2} cannot represent.  Whether its coefficient probabilities
are calibrated is a separate empirical question.  The benchmarks therefore
report failures as well as successes and do not treat the sampling-model choice
as proof of superiority over negative-binomial regression.
